\documentclass[11pt,a4paper]{article}

\usepackage[utf8]{inputenc}
\usepackage{amsmath,amssymb}
\usepackage[margin=2.5cm]{geometry}
\usepackage{graphicx}
\usepackage{xcolor}
\usepackage{enumitem}
\usepackage[skins]{tcolorbox}
\usepackage[colorlinks=true,linkcolor=violet!70!black,urlcolor=violet!70!black]{hyperref}

\definecolor{accent}{HTML}{534AB7}

\newtcolorbox{callout}{
  enhanced, frame hidden, boxrule=0pt, sharp corners,
  colback=accent!6, borderline west={2.5pt}{0pt}{accent},
  left=10pt, right=8pt, top=5pt, bottom=5pt, before skip=8pt, after skip=8pt
}

\setlist[itemize]{itemsep=2pt, topsep=4pt}
\setlist[enumerate]{itemsep=2pt, topsep=4pt}
\setlength{\parindent}{0pt}
\setlength{\parskip}{6pt}

\newcommand{\norm}[1]{\left\lVert #1 \right\rVert}

\title{Topic 11 --- SVD\\[2pt]\large (Singular Value Decomposition)}
\author{}
\date{}

\begin{document}

\maketitle
\vspace{-2.2em}

\begin{center}
  \textbf{Subject:} Linear Algebra for ML \quad|\quad
  \textbf{Status:} growing \quad|\quad
  \textbf{MOC:} Math
\end{center}

\hrule
\vspace{4pt}

\begin{abstract}
\noindent SVD decomposes any matrix (any shape) into three matrices: $U$, $\Sigma$, $V^T$.
It can also be written as a sum of rank-one layers sorted by importance.
The more layers we keep, the better the approximation of $X$.
\end{abstract}

% ------------------------------------------------------------
\section{SVD}

\begin{callout}
The process of \textbf{decomposing a matrix (any shape)} into the \textbf{product of three other matrices}.
\end{callout}

\[
  X = U \Sigma V^T
\]

\begin{itemize}
  \item $X = \begin{bmatrix} | & | & & | \\ x_1 & x_2 & \cdots & x_m \\ | & | & & | \end{bmatrix}$ = data matrix, with $n \gg m$
  \item $x_k \in \mathbb{R}^n$ (dimension)
  \item $U$, $V$ = \textbf{unitary matrices} (orthogonal matrices)
  \item $\Sigma$ = \textbf{diagonal matrix}
\end{itemize}

% ------------------------------------------------------------
\section{Expanded Form}

\[
  X =
  \begin{bmatrix} | & | & & | \\ u_1 & u_2 & \cdots & u_m \\ | & | & & | \end{bmatrix}
  \begin{bmatrix} \sigma_1 & & \\ & \ddots & \\ & & \sigma_m \\ & 0 & \end{bmatrix}
  \begin{bmatrix} v_1 & v_2 & \cdots & v_m \end{bmatrix}^T
\]

\begin{itemize}
  \item $U$ columns $\rightarrow$ \textbf{``eigen'' faces}
  \item $\Sigma$ $\rightarrow$ \textbf{ordered by importance}
  \item $V$ $\rightarrow$ \textbf{``eigen'' time series}
\end{itemize}

% ------------------------------------------------------------
\section{Unitary Matrices}

$U, V$ = unitary matrices

\[
  UU^T = U^TU = I_{n \times n} \qquad \text{(transpose)}
\]

\[
  VV^T = V^TV = I_{m \times m}
\]

\begin{callout}
$V$ has the same property as $U$.
\end{callout}

% ------------------------------------------------------------
\section{Naming the Matrices}

Singular values are ordered:

\[
  \sigma_1 \geq \sigma_2 \geq \sigma_3 \geq \dots \geq \sigma_m \geq 0
\]

So we call each of these matrices like so:

\[
  U \qquad \Sigma \qquad V^T
\]

\begin{itemize}
  \item $U$ = \textbf{left singular vectors}
  \item $\Sigma$ = \textbf{singular values}
  \item $V$ = \textbf{right singular vectors}
\end{itemize}

% ------------------------------------------------------------
\section{Geometric View}

\begin{figure}[ht]
  \centering
  \includegraphics[width=\linewidth]{svd-geometry-steps.png}
  \caption{SVD acts on the unit circle in three steps, shown for
  $A = \begin{bmatrix} 1.2 & 0.8 \\ 0.2 & 0.9 \end{bmatrix}$
  ($\sigma_1 \approx 1.61$, $\sigma_2 \approx 0.57$).}
\end{figure}

SVD = \textbf{rotate $\rightarrow$ stretch $\rightarrow$ rotate}, applied to the unit circle:

\begin{enumerate}
  \item $V^T$ $\rightarrow$ \textbf{rotate} ($v_1, v_2$ land on the axes)
  \item $\Sigma$ $\rightarrow$ \textbf{stretch} by $\sigma_1$ and $\sigma_2$ along the axes (circle $\rightarrow$ ellipse)
  \item $U$ $\rightarrow$ \textbf{rotate} again (the axes land on $u_1, u_2$)
\end{enumerate}

% ------------------------------------------------------------
\section{Interpretation of Components}

\begin{itemize}
  \item $U$ columns have the \textbf{same shape as $X$}, so they can be reshaped (e.g.\ back into images).
  \item $V^T$ rows $\rightarrow$ \textbf{mixtures of all values to take the $x$'s}:
\end{itemize}

\[
  V^T =
  \begin{bmatrix}
    \text{---} & v_1^T & \text{---} \\
    \text{---} & v_2^T & \text{---} \\
     & \vdots & \\
    \text{---} & v_m^T & \text{---}
  \end{bmatrix}
\]

% ------------------------------------------------------------
\section{Sum of Rank-One Matrices}

\[
  X = \sigma_1 u_1 v_1^T + \sigma_2 u_2 v_2^T + \dots + \sigma_m u_m v_m^T + 0
\]

\[
  X = U \Sigma V^T = \hat{U} \hat{\Sigma} V^T \quad \rightarrow \quad \text{Economy SVD}
\]

Each term is a column ($u_i$) times a row ($v_i^T$), scaled by $\sigma_i$:

\[
  X = \sigma_1 \underbrace{u_1}_{\text{column}} \underbrace{v_1^T}_{\text{row}}
  + \sigma_2 u_2 v_2^T + \dots
  + \underbrace{\sigma_r u_r v_r^T}_{\text{truncate at } r}
  + \dots + \sigma_m u_m v_m^T + 0
\]

\begin{itemize}
  \item These rank-one matrices are \textbf{layers that increasingly improve the approximation of $X$}.
  \item So each $\sigma_i u_i v_i^T$ is a \textbf{simple image}, and the layers are \textbf{sorted by importance} ($\sigma_1, \sigma_2, \dots$).
\end{itemize}

% ------------------------------------------------------------
\section{Why More ``Rank-One Matrices'' Gives a Better Image?}

\begin{itemize}
  \item \textbf{Layer 1} (biggest $\sigma$): captures the broad structure: overall light and dark regions.
  \item \textbf{Layers 2 to 20}: add the main shapes and edges.
  \item \textbf{Later layers}: add finer and finer detail.
  \item \textbf{The last layers}: have tiny $\sigma$ and mostly add noise and texture.
\end{itemize}

% ------------------------------------------------------------
\section{The Eckart--Young Theorem (1936)}

\begin{callout}
The \textbf{best approximation} to the matrix $X$ that has rank $r$.
\end{callout}

\[
  \underset{\tilde{X}\ \text{such that}\ \operatorname{rank}(\tilde{X}) = r}{\arg\min}
  \norm{X - \tilde{X}}_F = \tilde{U} \tilde{\Sigma} \tilde{V}^T
\]

\begin{itemize}
  \item $\arg\min$ $\rightarrow$ \textbf{smallest error}
  \item $F$ $\rightarrow$ \textbf{Frobenius norm}
\end{itemize}

$\rightarrow$ So the theorem \textbf{guarantees} that the best possible matrix approximation to $X$ of rank $r$ is the \textbf{first $r$ layers}.

% ------------------------------------------------------------
\section{Unitary Identity After Truncation}

\[
  \tilde{U}^T \tilde{U} = I_{r \times r}
\]

\[
  \tilde{U} \tilde{U}^T \neq I
\]

\begin{callout}
It's not an identity anymore.
\end{callout}

% ------------------------------------------------------------


\vspace{6pt}
\hrule
\vspace{6pt}



\end{document}